\documentclass[10pt,twocolumn]{article}
\usepackage{amsmath,fancyhdr,lipsum}
\pagestyle{fancy}\fancyhf{}\lhead{J. Chem. Phys. 160, 094101 (2024)}\rfoot{\thepage}
\title{Equivariant message passing potentials for biomolecular simulation}
\author{Ada Lovelace \and Alan Turing}
\begin{document}
\maketitle
\begin{abstract}
We introduce an E(3)-equivariant interatomic potential trained on density functional theory data for small proteins. It reaches chemical accuracy with a fraction of the training data required by invariant descriptors such as atom-centered symmetry functions.
\end{abstract}
\section{Introduction}
Classical force fields such as AMBER and CHARMM describe the potential energy surface with fixed functional forms. Machine-learned interatomic potentials instead learn the mapping from atomic configurations to energies and forces, following the high-dimensional neural network approach of Behler and Parrinello~\cite{bp}. \lipsum[1-2]
\section{Methods}
\subsection{Equivariant features}
The energy is a sum of atomic contributions,
\begin{equation}
E = \sum_i E_i(\{\mathbf{r}_{ij}\}_{j \in \mathcal{N}(i)}), \qquad \mathbf{F}_i = -\nabla_{\mathbf{r}_i} E,
\end{equation}
where features transform as irreducible representations of O(3). \lipsum[3-4]
\subsection{Training data}
We trained on 10\,000 configurations of alanine dipeptide sampled with the TIP3P water model. \lipsum[5-6]
\section{Results}
\lipsum[7-9]
\section{Conclusion}
\lipsum[10]
\section*{Acknowledgments}
We thank the funding agency for generous support of this work.
\begin{thebibliography}{9}
\bibitem{bp} J. Behler and M. Parrinello, Generalized neural-network representation of high-dimensional potential-energy surfaces, Phys. Rev. Lett. 98, 146401 (2007).
\bibitem{nequip} S. Batzner et al., E(3)-equivariant graph neural networks for data-efficient and accurate interatomic potentials, Nat. Commun. 13, 2453 (2022).
\bibitem{ff} J. A. Lemkul, The continuous evolution of biomolecular force fields, J. Chem. Theory Comput. (2024).
\end{thebibliography}
\appendix
\section{Hyperparameters}
All models used a radial cutoff of 5 \AA{} and 3 interaction layers. \lipsum[11]
\end{document}
